\documentclass{situationist-preprint}

\title{On Not Fearing Death}
\msid{EXI-05-a}
\msfamily{EXI-05}
\msdiscipline{Meaning, Death and the Absurd}
\mstype{Principal manuscript}
\mskeywords{fear of death; fittingness of emotions; regret; Epicurus; deprivation; uncertainty; Kagan}

\begin{document}
\maketitle

\begin{abstract}
Should we fear death? Epicurus held that we should not, because death is not bad for us. EXI-02 rejects that premise: death is bad for the one who dies when it deprives her of goods. It does not follow that fear of death is rational. We argue, using a fittingness account of emotions, that fear is fitting only towards a prospective evil whose occurrence is uncertain, and that one's own death as such, as distinct from particular threats to one's life, is certain. Fear of death is therefore not fitting, and so not rational, whatever one thinks of its badness. Regret at the prospect of death is a different attitude, a sorrow at goods one will not have, and its fittingness conditions are badness and prospect, not uncertainty. Since death is bad by deprivation, regret at the prospect of death may be rational, in proportion to what death will take. The result separates the Epicurean conclusion about fear from the Epicurean premise about badness.
\end{abstract}

\section{The problem}

Epicurus argued that it is foolish to be pained in prospect by what does not pain when it is present, and that death, which is nothing to us, should not be feared \citep[125]{EpicurusMen}. His argument for not fearing death runs through his claim that death is not bad for us. EXI-02 argues that this claim is false: death deprives the one who dies of the goods she would otherwise have had. If fear is the appropriate response to what is bad for one, it seems to follow that we should fear death after all. \citet{Kagan2012} argued that this does not follow, since fear has conditions of appropriateness beyond the badness of its object; \citet{Nussbaum1989} defended a qualified version of the Epicurean therapy. We argue for a precise version of the conclusion: fear of death is not rational, though regret at the prospect of death may be.

\section{Framework}

\begin{definition}[Fittingness]\label{d:fit}
An emotion is \emph{fitting} towards its object if the object has the evaluative features that the emotion represents it as having. An emotion is \emph{rational} only if it is fitting, given what the subject is in a position to know.
\end{definition}

\begin{principle}[Fear]\label{p:fear}
Fear represents its object as a prospective evil to oneself that may or may not come about. Fear of $X$ is fitting only if (i) $X$ is bad for the subject, (ii) $X$ is in the subject's future, and (iii) it is uncertain, for the subject, whether $X$ will occur.
\end{principle}

\begin{principle}[Regret]\label{p:regret}
Regret at a prospect represents its object as a loss of goods to oneself. Prospective regret about $X$ is fitting if $X$ is bad for the subject and in her future; it does not require uncertainty.
\end{principle}

Condition (iii) of Principle~\ref{p:fear} follows the analysis of fear in \citet{Gordon1987}, on which fear that $p$ requires that one not be certain that $p$: one fears that it will rain only while it is open whether it will. \citet{Kagan2012} lists the same condition among the conditions on appropriate fear, together with a nonnegligible probability and proportion. The condition explains the role of fear: it orients the subject to avoid or prepare for an evil that may still be averted or that may not come, and it is replaced, once the evil is certain, by dread, resignation or grief. Principle~\ref{p:regret} reflects the fact that one can be sad that something will happen when one knows that it will: one can regret, as it approaches, the end of a stay with friends.

\section{The argument}

\begin{lemma}\label{l:certain}
One's own death as such is certain.
\end{lemma}

\begin{proof}
Every human being is an organism, and every organism dies. No one is in a position to doubt that she will die, as distinct from doubting when or how.
\end{proof}

By \enquote{death} we mean, with EXI-02, the end of the psychological life that the organism sustains. REL-08 argues that survival of that end is coherent but that no evidence available before death could confirm it; the possibility of survival does not make the end itself uncertain, only what follows it.

\begin{theorem}\label{t:fear}
Fear of death as such is not fitting, and so not rational.
\end{theorem}

\begin{proof}
By Principle~\ref{p:fear}, fear of $X$ is fitting only if it is uncertain for the subject whether $X$ will occur. By Lemma~\ref{l:certain}, it is certain that one will die. So condition (iii) fails for fear of death as such, and by Definition~\ref{d:fit} fear of death is not rational.
\end{proof}

\begin{theorem}\label{t:regret}
Regret at the prospect of one's death may be rational, and it is fitting in proportion to the goods of which death will deprive one.
\end{theorem}

\begin{proof}
By EXI-02, one's death is bad for one when it deprives one of goods one would otherwise have had, and the degree of its badness is the value of what it takes. One's death is in one's future. So by Principle~\ref{p:regret}, prospective regret about it is fitting when death will deprive one of goods, and fitting to the degree that it does. Its fittingness does not depend on uncertainty, which Lemma~\ref{l:certain} excludes.
\end{proof}

\section{Objections and replies}

\begin{objection}[Premature death]\label{o:early}
Whether I will die soon is uncertain, and an early death would be a great deprivation. So fear of early death satisfies all the conditions of Principle~\ref{p:fear}.
\end{objection}

\begin{reply}
It does, and it is rational in proportion to the probability and the deprivation. But its object is not death as such; it is a particular threat, or the possibility that one's life will be short. Fear of a particular danger is a fear of a prospective evil whose occurrence is uncertain, and its rationality is the rationality of prudence. Theorem~\ref{t:fear} concerns the attitude the question is about: fear directed at one's mortality, at the fact that one will die. The distinction is the one that separates fearing a storm from fearing that the sun will set.
\end{reply}

\begin{objection}[Dread]\label{o:dread}
People dread certain evils, such as a scheduled operation. Dread is a form of fear, and it is rational when the evil is great.
\end{objection}

\begin{reply}
Dread of a scheduled operation is directed at what is uncertain about it: how much it will hurt, whether it will go well. Where nothing is uncertain, what remains is a sorrow at what is coming, which is the attitude of Principle~\ref{p:regret}. To call it fear is to name it by the agitation that often accompanies it, and Theorem~\ref{t:regret} grants that the sorrow is fitting.
\end{reply}

\begin{objection}[Asymmetry]\label{o:asym}
We do not regret having been born late. If regret at the prospect of death is rational, so should regret about our late birth be, by Lucretius's symmetry.
\end{objection}

\begin{reply}
Proposition~4.1 of EXI-02 shows that the alternatives on the two sides are not on a par: an earlier start is usually not a close possibility, and where it is, it is a different life whose value is not guaranteed to be greater. The bias towards the future \citep{BruecknerFischer1986} further explains why prospective regret is the more natural attitude, but the asymmetry of fitting regret does not depend on it.
\end{reply}

\section{Conclusion}

Fear of death is not rational, not because death is not bad, as Epicurus held, but because fear is fitting only towards an uncertain evil and death is certain. Regret at the prospect of death may be rational, since death is bad by deprivation, and is fitting in proportion to what death will take. Anxiety in the face of one's being towards death, in Heidegger's sense, is a different attitude that PHE-07 examines.

\begin{thebibliography}{99}\small
\bibitem[Brueckner and Fischer(1986)]{BruecknerFischer1986} Brueckner, A.~L. and Fischer, J.~M. (1986). Why is death bad? \emph{Philosophical Studies}, 50(2), 213--221.
\bibitem[Epicurus(Men.)]{EpicurusMen} Epicurus. \emph{Letter to Menoeceus}. In Diogenes Laertius, \emph{Lives of Eminent Philosophers} X.
\bibitem[Gordon(1987)]{Gordon1987} Gordon, R.~M. (1987). \emph{The Structure of Emotions: Investigations in Cognitive Philosophy}. Cambridge University Press.
\bibitem[Kagan(2012)]{Kagan2012} Kagan, S. (2012). \emph{Death}. Yale University Press.
\bibitem[Nussbaum(1989)]{Nussbaum1989} Nussbaum, M.~C. (1989). Mortal immortals: Lucretius on death and the voice of nature. \emph{Philosophy and Phenomenological Research}, 50(2), 303--351.
\end{thebibliography}

\end{document}
